%%%%%%%%%%%%%%%%%%%%%%%%%%%%%%%%%%%%%%%%%%%%%%%%%%%%%%%%%%%%%%%%%%%%%%%%%%%
%%  chapter2.tex  –  Chapter 2: Methods, Algorithms, and Mathematical
%%                   Frameworks for Robust Dynamic Graph Neural Networks
%%  v9 – Unified problem statement with progressive derivation.
%%       Decomposition into 7 sub-problems mapped to gaps.
%%       All proofs, algorithms, and complexity analyses preserved.
%%       Only \section headings go to TOC.
%%%%%%%%%%%%%%%%%%%%%%%%%%%%%%%%%%%%%%%%%%%%%%%%%%%%%%%%%%%%%%%%%%%%%%%%%%%

\chapter{Methods, Algorithms, and Mathematical Frameworks for Robust Dynamic Graph Neural Networks}
\label{ch:methods}

This chapter develops the complete theoretical core of the dissertation. It begins with a unified problem statement that is progressively derived from the fundamental graph classification problem by successively incorporating each of the four requirements (robustness, accuracy, efficiency, privacy) under the three operating conditions (adversarial environment, non-stationarity, distributed dynamics). The unified problem is then formally decomposed into seven sub-problems, each corresponding to one of the gaps identified in Chapter~\ref{ch:literature}. The remainder of the chapter presents the mathematical formulation, convergence and robustness proofs, algorithmic pseudocode, and complexity analyses for each of the seven methods that solve these sub-problems. The chapter concludes with a unified certification framework that couples all methods through consistency regularisation.


%%=====================================================================
\section{Unified Problem Statement}
\label{sec:problem_statement}
%%=====================================================================

The progressive derivation of the problem statement proceeds in seven stages, each introducing a new mathematical constraint that corresponds to a specific requirement or operating condition. The result is a single multi-objective optimisation problem whose decomposition yields the seven methods developed in this dissertation.


\subsection{Stage 1: Graph Classification on Dynamic Topologies}

Consider a temporal graph $\calG(t)=(\calV,\calE(t),\mathbf{X}(t))$ in which nodes $\calV$, edges $\calE(t)$, and features $\mathbf{X}(t)=\{\bx_v(t)\in\R^{d}\}_{v\in\calV}$ evolve over time. A graph neural network $f_\theta$ maps the graph at time $T$ to a prediction. The baseline classification problem is
\begin{equation}
\min_\theta\;
\E_{(\calG,y)\sim\calD}\!\bigl[
  \calL\!\bigl(f_\theta(\calG(T)),\,y\bigr)
\bigr],
\label{eq:ps_base}
\end{equation}
where $\calL$ is the cross-entropy loss and $\calD$ is the data-generating distribution. This formulation assumes that the graph is clean, static during inference, centrally available, and that the data distribution does not change. Each of these assumptions is violated in practice.


\subsection{Stage 2: Adversarial Robustness}

Under the adversarial environment (Condition C1), a deliberate opponent perturbs the graph. Incorporating Requirement R1 (robustness) converts the problem to a min-max formulation:
\begin{equation}
\min_\theta\;
\E_{(\calG,y)\sim\calD}\!\Bigl[
  \max_{\|\bdelta\|\le\epsilon_{\mathrm{adv}}}
    \calL\!\bigl(f_\theta(\calG(T)+\bdelta),\,y\bigr)
\Bigr],
\label{eq:ps_robust}
\end{equation}
where $\bdelta$ represents perturbations to node features, edge weights, or topological structure, and $\epsilon_{\mathrm{adv}}$ is the perturbation budget. This is the standard robust optimisation formulation~\cite{Madry2018,Goodfellow2015,Szegedy2014}.


\subsection{Stage 3: Continuous-Time Dynamics with Certified Bounds}

To obtain certified robustness rather than empirical adversarial training, the model's internal dynamics must be constrained. Requiring that $f_\theta$ evolves through a neural ODE with Lipschitz-bounded vector field adds the constraint:
\begin{equation}
\frac{d\bh_v(t)}{dt} = f_\theta\!\bigl(\bh_v(t),\,\mathbf{A}(t),\,t\bigr),
\qquad
\|f_\theta(\bh_1,t)-f_\theta(\bh_2,t)\| \le L_g\|\bh_1-\bh_2\|,
\label{eq:ps_ode}
\end{equation}
which, by Gr\"{o}nwall's inequality, guarantees $\|\bh_1(T)-\bh_2(T)\|\le e^{L_gT}\|\bx_1-\bx_2\|$. This transforms the uncertified adversarial training of Equation~\eqref{eq:ps_robust} into a certified robustness problem. \emph{Sub-problem 1 (Gap~1) is to solve Equations~\eqref{eq:ps_robust}--\eqref{eq:ps_ode} jointly, with multi-scale temporal dynamics and $O(1)$-memory adjoint training. This is addressed by Method M1: CT-TGNN.}


\subsection{Stage 4: Distributed Learning with Privacy}

Under distributed dynamics (Condition C3), the graph $\calG$ is partitioned across $K$ participants $\calD_1,\ldots,\calD_K$ who cannot share raw data. Incorporating Requirement R4 (privacy) requires that the training algorithm satisfies $(\epsilon,\delta)$-differential privacy while tolerating up to a fraction $q$ of Byzantine participants:
\begin{equation}
\min_\theta\;
\frac{1}{K}\sum_{k=1}^{K}\calL_k(\theta;\calD_k)
\quad\text{subject to}\quad
\Prob[\calM(\calD)\in S] \le e^{\epsilon}\Prob[\calM(\calD')\in S]+\delta,
\label{eq:ps_fed}
\end{equation}
where $\calM$ is the training mechanism and the constraint must hold for all neighbouring datasets $\calD,\calD'$. Additionally, the heterogeneous distributions across participants require alignment through the Wasserstein distance $W_1(\calD_k,\calD_{k'})$, which must itself be computed under the privacy constraint. \emph{Sub-problem 2 (Gap~2) is to solve Equation~\eqref{eq:ps_fed} with Byzantine-resilient aggregation and optimal-transport alignment while consuming the Lipschitz certificates from Sub-problem~1. This is addressed by Method M2: FedLLM-API.}


\subsection{Stage 5: Multi-Level Representation and Anti-Forgetting}

Incorporating Requirement R2 (accuracy) on temporal graphs with hierarchical structure requires learning representations at multiple semantic and temporal scales --- component, trace, and node levels --- while preventing catastrophic forgetting as the topology evolves:
\begin{equation}
\min_\theta\;
\sum_{g\in\{s,r,n\}}\lambda_g\,\calL_g(\theta;\calG_g)
+ \lambda_{\mathrm{ewc}}\sum_j F_j(\theta_j-\theta_{j}^{*})^2
+ \lambda_{\mathrm{con}}\!\sum_{(g,g')}\mathrm{KL}(p_g\|p_{g'}),
\label{eq:ps_triple}
\end{equation}
where $F_j$ is the Fisher information penalising departure from previously learned parameters $\theta^*$, and the KL-divergence term enforces cross-level consistency. \emph{Sub-problem 3 (Gap~3) is to solve Equation~\eqref{eq:ps_triple} with structured sparsity and quantisation-aware training. This is addressed by Method M3: TripleE-TGNN.}


\subsection{Stage 6: Computational Efficiency}

Requirement R3 (efficiency) demands that inference complexity is at most $\calO(L\log L)$ rather than the $\calO(L^2)$ of transformer-based models. This is achieved by reformulating the sequential processing as a selective state-space model with input-dependent transitions:
\begin{equation}
\bh_{k} = \bar{\bA}(u_k)\,\bh_{k-1} + \bar{\bB}(u_k)\,u_k,
\qquad
\bar{\bA}(u_k)=\exp(\bA(u_k)\Delta),
\label{eq:ps_ssm}
\end{equation}
whose convolutional kernel is evaluated through the Fast Fourier Transform. This must be combined with progressive adversarial robustness distillation and PAC-Bayesian generalisation certificates. \emph{Sub-problem 4 (Gap~4) is to solve the efficient inference problem of Equation~\eqref{eq:ps_ssm} with formal robustness guarantees under training-time poisoning. This is addressed by Method M4: MambaShield.}


\subsection{Stage 7: Non-Stationary Adaptation with Uncertainty}

Under non-stationarity (Condition C2), the distribution $\calD_t$ changes over time. The model must adapt online while quantifying both epistemic and aleatoric uncertainty and achieving sub-linear regret:
\begin{equation}
R(T) = \max_{f^*}\sum_{t=1}^{T}\calL(f_{\theta_t},y_t) - \sum_{t=1}^{T}\calL(f^*,y_t) \le \calO(\sqrt{T\log|\calS|}),
\label{eq:ps_regret}
\end{equation}
where the comparison is against the best fixed policy in hindsight. The predictive uncertainty decomposes as $U_{\mathrm{total}} = U_{\mathrm{epi}} + U_{\mathrm{ale}}$, with epistemic uncertainty driving drift detection and aleatoric uncertainty capturing inherent noise.

\emph{Sub-problem 5 (Gap~5) concerns forward compatibility under protocol regime change and is addressed by a domain-specific foundation model (CyberSecLLM). Sub-problem 6 (Gap~6) is to solve the online adaptation problem of Equation~\eqref{eq:ps_regret} with integrated epistemic--aleatoric decomposition and Bayesian posterior updates. This is addressed by Methods M5: Stochastic Transformer and M6: UC-HGP.}


\subsection{The Complete Unified Problem}

Combining all stages yields the unified multi-objective optimisation problem of this dissertation:

\begin{equation}
\boxed{
\begin{aligned}
\min_\theta\;&
\underbrace{\E_{(\calG,y)\sim\calD_t}\!\Bigl[
  \max_{\|\bdelta\|\le\epsilon}
    \calL\!\bigl(f_\theta(\calG+\bdelta),y\bigr)
\Bigr]}_{\text{Adversarial robustness (R1, C1)}}
+ \lambda_{\mathrm{acc}}\underbrace{\sum_g\calL_g(\theta;\calG_g)}_{\text{Multi-level accuracy (R2)}}
\\[6pt]
&+ \lambda_{\mathrm{eff}}\underbrace{\mathrm{Complexity}(f_\theta)}_{\text{Efficiency $\le\calO(L\log L)$ (R3)}}
+ \lambda_{\mathrm{priv}}\underbrace{\mathrm{PrivCost}(\epsilon_{\mathrm{DP}},\delta)}_{\text{Differential privacy (R4, C3)}}
\\[6pt]
&+ \lambda_{\mathrm{drift}}\underbrace{R(T)}_{\text{Non-stationary regret (C2)}}
+ \lambda_{\mathrm{ewc}}\underbrace{\sum_j F_j(\theta_j-\theta_j^*)^2}_{\text{Anti-forgetting}}
+ \lambda_{\mathrm{con}}\underbrace{\sum_{i<j}\calL_{ij}(\theta_i,\theta_j)}_{\text{Method coupling (M7)}}
\end{aligned}
}
\label{eq:ps_unified}
\end{equation}

\noindent subject to the constraints:
\begin{itemize}
\item Lipschitz-bounded ODE dynamics: $\|f_\theta(\bh_1,t)-f_\theta(\bh_2,t)\|\le L_g\|\bh_1-\bh_2\|$ (certified robustness)
\item Differential privacy: $\Prob[\calM(\calD)\in S]\le e^{\epsilon}\Prob[\calM(\calD')\in S]+\delta$ (privacy)
\item Byzantine resilience: convergence under fraction $q<1/2$ corrupted participants
\item Sub-linear regret: $R(T)\le\calO(\sqrt{T\log|\calS|})$ (adaptation)
\end{itemize}

\noindent\textbf{Decomposition.} This unified problem is intractable as a monolithic optimisation. However, it admits a natural decomposition into seven sub-problems, one per identified gap, that can be solved by alternating optimisation with consistency coupling. Table~\ref{tab:decomposition} maps each sub-problem to its gap, the terms in Equation~\eqref{eq:ps_unified} it addresses, and the method that solves it.

\begin{table}[htbp]
\centering
\caption{Decomposition of the unified problem statement into seven sub-problems.}
\label{tab:decomposition}
\small
\begin{tabular}{p{0.5cm} p{3.0cm} p{4.5cm} p{5.0cm}}
\toprule
\textbf{SP} & \textbf{Method} & \textbf{Terms Addressed} & \textbf{Key Constraint} \\
\midrule
1 & M1: CT-TGNN & Adversarial robustness + ODE dynamics & Lipschitz certificate: $e^{L_gT}\|\delta\bx\|$ \\
2 & M2: FedLLM-API & Privacy cost + Byzantine aggregation & $(\epsilon,\delta)$-DP + OT alignment \\
3 & M3: TripleE-TGNN & Multi-level accuracy + anti-forgetting & EWC Fisher penalty + cross-level KL \\
4 & M4: MambaShield & Efficiency + adversarial robustness & $\calO(L\log L)$ + PAC-Bayes bound \\
5 & CyberSecLLM & Forward compatibility (domain-specific) & Multi-task ODE-S6 blocks \\
6 & M5/M6 & Non-stationary regret + uncertainty & Sub-linear regret + epi/ale decomposition \\
7 & M7: Certification & Method coupling + game-theoretic & Stackelberg equilibrium + consistency \\
\bottomrule
\end{tabular}
\end{table}

\emph{Sub-problem 7 (Gap~7) is to couple all methods through the consistency regularisation term $\sum_{i<j}\calL_{ij}$ in Equation~\eqref{eq:ps_unified} and prove convergence of the alternating optimisation. This is addressed by Method M7: Certification and the unified framework developed in Section~\ref{sec:unified}.}

The remainder of this chapter develops each method in the order of the sub-problem chain: M1 (Section~\ref{sec:ct_tgnn}), M2 (Section~\ref{sec:fedllm}), M3 (Section~\ref{sec:triple_e}), M4 (Section~\ref{sec:mamba}), M5/M6 (Section~\ref{sec:stochastic}), forward-compatible detection (Section~\ref{sec:pq_idps}), and M7 with the unified framework (Section~\ref{sec:game}).



%%=====================================================================
\section{Method M1: CT-TGNN --- Continuous-Time Temporal Graph Neural Dynamics}
\label{sec:ct_tgnn}
%%=====================================================================

This section solves Sub-problem~1 by developing a graph neural network whose node representations evolve according to a neural ordinary differential equation, yielding the certified robustness bound from the Lipschitz constraint in Equation~\eqref{eq:ps_ode}. The method addresses Gap~1 (Adversarial $\times$ Robustness; Non-stationary $\times$ Robustness).

\end{equation}
where $\bh_v(t)\in\R^{d_h}$ is the embedding of node~$v$ at time~$t$, $\calN_v(t)$ denotes the neighbourhood of~$v$ under the current topology, $\mathbf{A}(t)$ encodes the adjacency structure, and $f_\theta$ is a neural network parameterised by~$\theta$.

For heterogeneous graphs the right-hand side employs type-specific message passing. Let $v$ belong to type~$k$ and let $\calN_v^{(k')}(t)$ denote its neighbours of type~$k'$. The dynamics are
\begin{equation}
\frac{d\bh_v^{(k)}(t)}{dt}
= \sum_{k'=1}^{K}\sum_{u\in\calN_v^{(k')}(t)}
  \alpha_{uv}^{(k,k')}(t)\;\mathrm{MSG}_{k,k'}\!\bigl(
    \bh_u^{(k')}(t),\;\mathbf{e}_{uv}(t)\bigr),
\label{eq:ct_tgnn_hetero}
\end{equation}
where the attention coefficients~\cite{Velickovic2018}
\begin{equation}
\alpha_{uv}^{(k,k')}(t)
= \frac{
  \exp\!\bigl(\mathrm{LeakyReLU}\!\bigl(
    \mathbf{a}_{k,k'}^\top
    [\bW_k\bh_v^{(k)}(t)\|\bW_{k'}\bh_u^{(k')}(t)]
  \bigr)\bigr)
}{
  \sum_{w\in\calN_v^{(k')}(t)}
  \exp\!\bigl(\mathrm{LeakyReLU}\!\bigl(
    \mathbf{a}_{k,k'}^\top
    [\bW_k\bh_v^{(k)}(t)\|\bW_{k'}\bh_w^{(k')}(t)]
  \bigr)\bigr)
}
\label{eq:ct_tgnn_attn}
\end{equation}
weight neighbour importance adaptively, $\bW_k,\bW_{k'}$ are learnable projections, and $\|$ denotes concatenation.

\subsection{Multi-Scale Temporal Dynamics}

Phenomena that unfold across multiple time scales, from sub-second transients to month-long trends, require parallel ODE branches with distinct time constants. For scales $s=1,\dots,S$ with constants $\tau_1<\tau_2<\cdots<\tau_S$ the dynamics are
\begin{equation}
\frac{d\bh_v^{(s)}(t)}{dt}
= \frac{1}{\tau_s}\;f_{\theta_s}\!\bigl(
  \bh_v^{(s)}(t),\;\{\bh_u^{(s)}(t)\}_{u\in\calN_v(t)},\;
  \mathbf{A}(t),\;t\bigr),
\label{eq:ct_tgnn_multi}
\end{equation}
and the final embedding aggregates across scales through learned attention,
\begin{equation}
\bh_v(t)
= \sum_{s=1}^{S}\beta_s(t)\;\bh_v^{(s)}(t),
\qquad
\beta_s(t) = \mathrm{Softmax}_s\!\bigl(\mathbf{w}_s^\top\bh_v^{(s)}(t)\bigr),
\label{eq:ct_tgnn_scale_fuse}
\end{equation}
so that the relative contribution of each time scale adapts to the local dynamics.

\subsection{Adjoint Sensitivity Method for Memory-Efficient Training}

Training requires gradients of a loss $\calL=\calL(\bh(T))$ with respect to~$\theta$ through the ODE solution. Direct back-propagation through the numerical integrator stores all intermediate states, which is prohibitive for long horizons. The adjoint method~\cite{Pontryagin1962,Chen2018} defines an adjoint state $\mathbf{a}(t)=\partial\calL/\partial\bh(t)$ satisfying the backward ODE
\begin{equation}
\frac{d\mathbf{a}(t)}{dt}
= -\mathbf{a}(t)^\top\frac{\partial f_\theta}{\partial\bh}(\bh(t),t),
\label{eq:adjoint_ode}
\end{equation}
integrated from $\mathbf{a}(T)=\partial\calL/\partial\bh(T)$ to $\mathbf{a}(0)$. Parameter gradients follow from
\begin{equation}
\frac{\partial\calL}{\partial\theta}
= -\int_0^T \mathbf{a}(t)^\top
  \frac{\partial f_\theta}{\partial\theta}(\bh(t),t)\;dt.
\label{eq:adjoint_param}
\end{equation}
Memory cost is proportional to the model size rather than the trajectory length, enabling training on arbitrarily long temporal horizons.

\subsection{Temporal Adaptive Batch Normalisation}

Standard batch normalisation computes statistics over a fixed mini-batch, which is incompatible with continuous-time dynamics in which the distribution of hidden states changes continuously. Temporal adaptive batch normalisation computes running statistics over a sliding temporal window $[t-\Delta,t]$,
\begin{equation}
\mathrm{TABN}\!\bigl(\bh(t)\bigr)
= \gamma(t)\;\frac{\bh(t)-\mu_{[t-\Delta,t]}}
  {\sqrt{\sigma^2_{[t-\Delta,t]}+\varepsilon}}
+ \beta(t),
\label{eq:tabn}
\end{equation}
where $\mu_{[t-\Delta,t]}$ and $\sigma^2_{[t-\Delta,t]}$ are the windowed mean and variance, and the affine parameters $\gamma(t),\beta(t)$ themselves evolve through auxiliary ODEs $d\gamma/dt=g_\gamma(\bh(t))$, $d\beta/dt=g_\beta(\bh(t))$ driven by learned adaptation networks.

\subsection{Convergence and Robustness Analysis}

\begin{theorem}[CT-TGNN Convergence]
\label{thm:ct_tgnn_conv}
Suppose the dynamics function $f_\theta$ is $L_f$-Lipschitz in~$\bh$ and $L_\theta$-Lipschitz in~$\theta$. Under gradient descent with learning rate $\eta<2/(L_f L_\theta)$, the training loss converges to a stationary point at rate $\calO(1/\sqrt{T})$, where $T$ denotes the number of gradient steps.
\end{theorem}

\noindent\textit{Proof sketch.} The loss $\calL(\theta)$ is evaluated through the ODE solution map $\bh(T;\theta)$. By the chain rule and the Lipschitz properties of~$f_\theta$, the gradient $\nabla_\theta\calL$ is bounded and Lipschitz-continuous, placing the problem in the standard framework of smooth non-convex optimisation. Applying the descent lemma with the stated learning rate yields the $\calO(1/\sqrt{T})$ rate to an $\epsilon$-stationary point. The full proof appears in Appendix~A.\qed

\begin{theorem}[Lipschitz Robustness Certificate]
\label{thm:ct_tgnn_rob}
Let $f_\theta$ be $L_g$-Lipschitz in its first argument and let the integration horizon be~$T$. For any two inputs $\bx_1,\bx_2$ the corresponding embeddings satisfy
\begin{equation}
\|\bh_1(T)-\bh_2(T)\|
\le e^{L_g T}\;\|\bx_1-\bx_2\|.
\label{eq:lip_cert}
\end{equation}
Consequently, if the classifier head has Lipschitz constant~$L_c$, the end-to-end model is $L_c\,e^{L_g T}$-Lipschitz, and the prediction is certified constant within a ball of radius $\rho = \Delta_{\min}/(L_c\,e^{L_g T})$, where $\Delta_{\min}$ is the margin between the top two logits.
\end{theorem}

\noindent\textit{Proof sketch.} Let $\bm{\eta}(t)=\bh_1(t)-\bh_2(t)$. Then $d\bm{\eta}/dt = f_\theta(\bh_1,t)-f_\theta(\bh_2,t)$, so $d\|\bm{\eta}\|/dt\le L_g\|\bm{\eta}\|$ by the Lipschitz property. Gr\"{o}nwall's inequality gives $\|\bm{\eta}(T)\|\le e^{L_g T}\|\bm{\eta}(0)\|$, and $\bm{\eta}(0)=\bx_1-\bx_2$. Composing with the classifier head completes the certificate. The full derivation with spectral-norm constraints enforcing the Lipschitz bound during training appears in Appendix~A.\qed

\subsection{Computational Complexity}

Solving the graph ODE with an adaptive Runge-Kutta integrator at tolerance~$\varepsilon_{\mathrm{ode}}$ requires $\calO(\varepsilon_{\mathrm{ode}}^{-1/5})$ function evaluations. Each evaluation performs message passing at cost $\calO(|\calV|\,d\,d_h)$ for bounded degree~$d$ and embedding dimension~$d_h$. The total training complexity per epoch is therefore
\begin{equation}
\calO\!\bigl(|\calV|\,d\,d_h\;\varepsilon_{\mathrm{ode}}^{-1/5}\;T_{\mathrm{train}}\bigr),
\label{eq:ct_tgnn_complex}
\end{equation}
which is near-linear in graph size when $|\calE|=\calO(|\calV|)$.

\begin{algorithm}[htbp]
\caption{CT-TGNN Training with Adjoint Sensitivity}\label{alg:ct_tgnn}
\KwInput{Temporal graph sequence $\{\calG(t_i)\}_{i=1}^N$, labels $\{y_i\}$, learning rate $\eta$, integration horizon $T$, tolerance $\varepsilon_{\mathrm{ode}}$}
\KwOutput{Trained parameters $\theta^*$}
Initialise parameters $\theta_0$ and multi-scale time constants $\tau_1<\cdots<\tau_S$\;
\For{epoch $= 1,2,\ldots,E$}{
  \For{each mini-batch $\{(\calG(t_j),y_j)\}$}{
    Set initial embeddings $\bh_v(0) \leftarrow \mathrm{Embed}(\bx_v(t_j))$ for all $v\in\calV$\;
    \For{scale $s=1,\ldots,S$}{
      Solve $d\bh_v^{(s)}/dt = \tau_s^{-1}f_{\theta_s}(\bh_v^{(s)},\calN_v,\bA,t)$ from $0$ to $T$ using adaptive RK with tolerance $\varepsilon_{\mathrm{ode}}$\;
    }
    Compute fused embeddings $\bh_v(T) = \sum_s \beta_s(T)\,\bh_v^{(s)}(T)$\;
    Compute classification loss $\calL = \calL_{\mathrm{CE}}(f_{\mathrm{head}}(\bh(T)),\,y)$\;
    Set adjoint terminal condition $\mathbf{a}(T) = \partial\calL/\partial\bh(T)$\;
    Solve adjoint ODE backward: $d\mathbf{a}/dt = -\mathbf{a}^\top(\partial f_\theta/\partial\bh)$ from $T$ to $0$\;
    Accumulate parameter gradients $\partial\calL/\partial\theta = -\int_0^T \mathbf{a}^\top(\partial f_\theta/\partial\theta)\,dt$\;
    Update $\theta \leftarrow \theta - \eta\,\partial\calL/\partial\theta$\;
  }
}
Compute Lipschitz certificate $\rho = \Delta_{\min}/(L_c\,e^{L_g T})$ for each test input\;
\Return $\theta^*,\;\{\rho_i\}$
\end{algorithm}


%%=====================================================================
\section{Method M2: FedLLM-API --- Federated Graph Optimisation with Privacy}
\label{sec:fedllm}
%%=====================================================================

This section solves Sub-problem~2 by developing a federated learning algorithm that achieves the privacy constraint and Byzantine resilience of Equation~\eqref{eq:ps_fed} while consuming the Lipschitz certificates produced by M1. The method addresses Gap~2 (Distributed $\times$ Robustness; Distributed $\times$ Privacy; Distributed $\times$ Efficiency).

\subsection{Byzantine-Robust Aggregation}

Consider $K$ participants with local datasets $\calD_1,\dots,\calD_K$, of which at most a fraction~$q$ are Byzantine. At communication round~$t$ the server broadcasts global model~$\bw_t$ and each honest client~$k$ computes a local update through gradient descent on its own data, $\bw_{t+1}^{(k)}=\bw_t - \eta\nabla_{\bw}\calL_k(\bw_t;\calD_k)$. Malicious clients submit arbitrary vectors. The server aggregates through coordinate-wise trimmed mean~\cite{Yin2018,Blanchard2017,Li2020fedprox}: for each parameter dimension~$j$ the client values are sorted and the $b=\lfloor Kq\rfloor$ largest and smallest are discarded, yielding
\begin{equation}
\bar{w}_{t+1,j}
= \frac{1}{K-2b}\sum_{i=b+1}^{K-b}w_{t+1,j}^{(i)},
\label{eq:fed_trim}
\end{equation}
where $w_{t+1,j}^{(i)}$ is the $i$-th order statistic.

\begin{theorem}[Byzantine-Robust Convergence]
\label{thm:fed_conv}
Assume the global loss $F(\bw)=\frac{1}{K}\sum_{k=1}^K F_k(\bw)$ is $\mu$-strongly convex and $L$-smooth, and the fraction of Byzantine clients satisfies $f<q<1/2$. With learning rate $\eta\le 1/L$, coordinate-wise trimmed mean achieves
\begin{equation}
\E\!\bigl[\|\bw_T-\bw^*\|^2\bigr]
\le \Bigl(1-\frac{\mu\eta}{2}\Bigr)^{\!T}\|\bw_0-\bw^*\|^2
  + \frac{2\eta}{\mu(K{-}2b)}
    \Bigl(\sum_{k=1}^{K-2b}\sigma_k^2 + C_{\mathrm{byz}}\Bigr),
\label{eq:fed_conv_bound}
\end{equation}
where $\sigma_k^2$ is the gradient variance of honest client~$k$ and $C_{\mathrm{byz}}$ bounds the Byzantine influence.
\end{theorem}

\noindent\textit{Proof sketch.} The key step is showing that the trimmed mean of honest-plus-Byzantine gradient vectors is close to the true mean of honest gradients. Let $\bar{\mathbf{g}}_H$ be the mean of honest gradients and $\bar{\mathbf{g}}_{\mathrm{TM}}$ the trimmed-mean output. Because at most $b$ of the $2b$ discarded entries can be honest, the error $\|\bar{\mathbf{g}}_{\mathrm{TM}}-\bar{\mathbf{g}}_H\|$ is bounded by the order statistics of honest gradient norms via concentration inequalities. Substituting into the standard strongly convex descent recursion yields the stated bound. Full details are in Appendix~A.\qed

\subsection{Differential Privacy Through Calibrated Noise}

Each client clips its gradient norm to a threshold~$C$ and adds Gaussian noise~\cite{Abadi2016,Dwork2006,Dwork2014} before transmission,
\begin{equation}
\bar{\mathbf{g}}_k = \mathbf{g}_k\Big/\max\!\Bigl(1,\frac{\|\mathbf{g}_k\|}
{C}\Bigr),
\qquad
\tilde{\mathbf{g}}_k = \bar{\mathbf{g}}_k + \calN(\mathbf{0},\sigma^2 C^2\mathbf{I}),
\label{eq:fed_dp_noise}
\end{equation}
where the noise scale~$\sigma$ is determined by the target privacy parameters $(\epsilon,\delta)$. Privacy amplification by sub-sampling with batch probability~$q_b$ and composition over $T$ rounds through the moments accountant~\cite{Mironov2017,Balle2018} yields the overall privacy guarantee
\begin{equation}
\epsilon_{\mathrm{total}}
= \min_{\alpha>1}\;
  \frac{1}{\alpha-1}\log\frac{1}{\delta}
  + \frac{T\,q_b^2\,\alpha}{2\sigma^2}.
\label{eq:fed_privacy}
\end{equation}

\subsection{Optimal Transport for Distribution Alignment}

Heterogeneous client distributions degrade convergence. Optimal transport aligns them by computing the Wasserstein barycenter. Given discrete source samples $\{\bx_i\}_{i=1}^{n_s}$ and target samples $\{\by_j\}_{j=1}^{n_t}$, the discrete transport plan solves
\begin{equation}
\mathbf{T}^* = \arg\min_{\mathbf{T}\in\R_+^{n_s\times n_t}}
\langle\mathbf{C},\mathbf{T}\rangle_F
\quad\text{s.t.}\quad
\mathbf{T}\mathbf{1}_{n_t}=\mathbf{a},\;
\mathbf{T}^\top\mathbf{1}_{n_s}=\mathbf{b},
\label{eq:fed_ot}
\end{equation}
where $C_{ij}=c(\bx_i,\by_j)$ is the cost matrix and $\mathbf{a},\mathbf{b}$ lie in probability simplices. The Sinkhorn algorithm~\cite{Cuturi2013,Chizat2020,Villani2009} computes the entropically regularised solution through alternating scaling,
\begin{equation}
\mathbf{u}^{(\ell+1)} = \mathbf{a}\oslash(\mathbf{K}\mathbf{v}^{(\ell)}),
\qquad
\mathbf{v}^{(\ell+1)} = \mathbf{b}\oslash(\mathbf{K}^\top\mathbf{u}^{(\ell+1)}),
\label{eq:sinkhorn}
\end{equation}
where $\mathbf{K}=\exp(-\mathbf{C}/\lambda)$ is the entropic kernel, $\oslash$ denotes element-wise division, and $\lambda>0$ controls regularisation. After convergence the transport plan is $\mathbf{T}^*=\mathrm{diag}(\mathbf{u})\mathbf{K}\,\mathrm{diag}(\mathbf{v})$.

\subsection{Differentially Private Optimal Transport}

To protect sample locations, Gaussian noise is added to the empirical distribution before transport computation,
\begin{equation}
\tilde{\mu}_s = \frac{1}{n_s}\sum_{i=1}^{n_s}\delta_{\bx_i+\bm{\xi}_i},
\qquad \bm{\xi}_i\sim\calN(\mathbf{0},\sigma_{\mathrm{DP}}^2\mathbf{I}),
\label{eq:fed_dp_ot_dist}
\end{equation}
and Laplace noise is added to the resulting transport plan,
\begin{equation}
\tilde{\mathbf{T}} = \mathbf{T}^*
+ \mathrm{Lap}\!\Bigl(\frac{\Delta_T}{\epsilon_{\mathrm{OT}}}\Bigr),
\label{eq:fed_dp_ot_plan}
\end{equation}
where the sensitivity $\Delta_T=\max_{\calD,\calD'}\|\mathbf{T}^*(\calD)-\mathbf{T}^*(\calD')\|_1$ bounds the maximum change under single-record substitution and $\epsilon_{\mathrm{OT}}$ is the privacy budget allocated to transport.

\subsection{Large Language Model Integration for Zero-Shot Generalisation}

Semantic reasoning about novel categories is achieved by constructing a natural language prompt from the structural features of an observation and querying a large language model to produce a probability $p_{\mathrm{LLM}}(\mathrm{class}\mid\mathbf{r})$. The final prediction is the ensemble
\begin{equation}
p_{\mathrm{final}}
= \lambda_{\mathrm{LLM}}\;p_{\mathrm{LLM}}(\mathrm{class}\mid\mathbf{r})
  + (1-\lambda_{\mathrm{LLM}})\;p_{\mathrm{GNN}}(\mathrm{class}\mid\calG(\mathbf{r})),
\label{eq:fed_llm}
\end{equation}
where $\calG(\mathbf{r})$ is the graph constructed from the observation and $\lambda_{\mathrm{LLM}}$ balances the two sources.

\begin{algorithm}[htbp]
\caption{FedLLM-API: Byzantine-Robust Federated Optimisation with DP-OT}\label{alg:fedllm}
\KwInput{Clients $\calD_1,\ldots,\calD_K$, fraction Byzantine $q$, clipping threshold $C$, privacy budget $(\epsilon,\delta)$, Sinkhorn regularisation $\lambda$, communication rounds $T$}
\KwOutput{Global model $\bw_T$}
Initialise global model $\bw_0$; compute noise scale $\sigma = C\sqrt{2\ln(1.25/\delta)}/\epsilon$\;
\For{round $t = 0,1,\ldots,T{-}1$}{
  Broadcast $\bw_t$ to all clients\;
  \For{each client $k = 1,\ldots,K$ \textup{(in parallel)}}{
    Compute local gradient $\mathbf{g}_k = \nabla_{\bw}\calL_k(\bw_t;\calD_k)$\;
    Clip: $\bar{\mathbf{g}}_k \leftarrow \mathbf{g}_k / \max(1,\|\mathbf{g}_k\|/C)$\;
    Add noise: $\tilde{\mathbf{g}}_k \leftarrow \bar{\mathbf{g}}_k + \calN(\mathbf{0},\sigma^2 C^2\mathbf{I})$\;
    Send $\tilde{\mathbf{g}}_k$ to server\;
  }
  Set $b \leftarrow \lfloor Kq\rfloor$\;
  \For{each dimension $j$}{
    Sort $\{\tilde{g}_{k,j}\}_{k=1}^K$; compute trimmed mean $\bar{g}_{t,j} = \frac{1}{K-2b}\sum_{i=b+1}^{K-b}\tilde{g}_{j}^{(i)}$\;
  }
  \If{heterogeneity detected via $W_1$ estimate}{
    Compute noisy source distributions $\tilde{\mu}_k$ by adding $\calN(\mathbf{0},\sigma_{\mathrm{DP}}^2\mathbf{I})$ to samples\;
    Run Sinkhorn iterations to compute regularised transport plan $\mathbf{T}^*$\;
    Add Laplace noise: $\tilde{\mathbf{T}} \leftarrow \mathbf{T}^* + \mathrm{Lap}(\Delta_T/\epsilon_{\mathrm{OT}})$\;
    Apply alignment correction to $\bar{\mathbf{g}}_t$ using $\tilde{\mathbf{T}}$\;
  }
  Update global model: $\bw_{t+1} \leftarrow \bw_t - \eta\,\bar{\mathbf{g}}_t$\;
}
\Return $\bw_T$
\end{algorithm}


%%=====================================================================
\section{Method M3: TripleE-TGNN --- Multi-Level Temporal Graph Embeddings}
\label{sec:triple_e}
%%=====================================================================

This section solves Sub-problem~3 by developing multi-level graph embeddings that address the accuracy and anti-forgetting terms in Equation~\eqref{eq:ps_triple}. The method addresses Gap~3 (Non-stationary $\times$ Accuracy; Non-stationary $\times$ Efficiency).

\subsection{Hierarchical Graph Representation}

Many systems exhibit a natural hierarchy: a coarse level of functional components, an intermediate level of operational traces that span components, and a fine level of infrastructure hosts. This structure is captured by three coupled temporal graphs: a component graph $\calG_s(t)=(\calV_s,\calE_s(t),\mathbf{X}_s(t))$, a trace graph $\calG_r(t)=(\calV_r,\calE_r(t),\mathbf{X}_r(t))$, and a host graph $\calG_n(t)=(\calV_n,\calE_n(t),\mathbf{X}_n(t))$. Bipartite coupling graphs $\calB_{sr},\calB_{sn},\calB_{rn}$ link entities across levels: a component deploys on a host, a trace traverses a component, and a trace executes on a host.

\subsection{Multi-Level Embedding Learning}

Each granularity $g\in\{s,r,n\}$ maintains a dedicated encoder,
\begin{equation}
\bh_v^{(g)}(t)
= \mathrm{ENC}_g\!\bigl(
  \bx_v^{(g)}(t),\;
  \{\bh_u^{(g)}(t{-}1)\}_{u\in\calN_v^{(g)}(t)},\;
  \Delta t
\bigr),
\label{eq:triple_enc}
\end{equation}
where $\mathrm{ENC}_g$ is a granularity-specific temporal graph neural network. Cross-granularity message passing integrates information across abstraction levels through bipartite attention,
\begin{equation}
\mathbf{m}_v^{(g\to g')}(t)
= \sum_{u\in\calN_v^{(g,g')}}
  \alpha_{uv}^{(g,g')}(t)\;\bW_{g\to g'}\bh_u^{(g)}(t),
\label{eq:triple_cross_msg}
\end{equation}
with attention scores
\begin{equation}
\alpha_{uv}^{(g,g')}(t)
= \mathrm{Softmax}_u\!\Biggl(
  \frac{(\mathbf{Q}_{g'}\bh_v^{(g')}(t))^\top
        (\mathbf{K}_g\bh_u^{(g)}(t))}
       {\sqrt{d_k}}
\Biggr),
\label{eq:triple_cross_attn}
\end{equation}
where $\mathbf{Q}_{g'},\mathbf{K}_g$ are query and key projections and $\calN_v^{(g,g')}$ denotes neighbours in the coupling graph. The final embedding fuses intra- and cross-granularity information through layer normalisation,
\begin{equation}
\bh_v^{(g)}(t)
\leftarrow
\mathrm{LayerNorm}\!\Bigl(
  \bh_v^{(g)}(t)
  + \sum_{g'\neq g}\mathbf{m}_v^{(g'\to g)}(t)
\Bigr).
\label{eq:triple_fuse}
\end{equation}

\subsection{Dual Temporal-Spatial Attention}

Within each granularity, temporal attention aggregates historical states over a window of width~$w$,
\begin{equation}
\tilde{\bh}_v^{(g)}(t)
= \sum_{\tau=t-w}^{t-1}
  \alpha_\tau^{\mathrm{temp}}\;\bh_v^{(g)}(\tau),
\qquad
\alpha_\tau^{\mathrm{temp}}
= \mathrm{Softmax}_\tau\!\bigl(
  \mathbf{w}_{\mathrm{temp}}^\top\bh_v^{(g)}(\tau)\bigr),
\label{eq:triple_temp_attn}
\end{equation}
while spatial attention aggregates the current neighbourhood,
\begin{equation}
\hat{\bh}_v^{(g)}(t)
= \sum_{u\in\calN_v^{(g)}(t)}
  \alpha_u^{\mathrm{spat}}\;\bh_u^{(g)}(t),
\qquad
\alpha_u^{\mathrm{spat}}
= \mathrm{Softmax}_u\!\bigl(
  \mathbf{w}_{\mathrm{spat}}^\top
  [\bh_v^{(g)}(t)\|\bh_u^{(g)}(t)]\bigr).
\label{eq:triple_spat_attn}
\end{equation}
The combined representation is
\begin{equation}
\bh_v^{(g)}(t)
= \mathrm{MLP}\!\bigl(
  [\tilde{\bh}_v^{(g)}(t)\|
   \hat{\bh}_v^{(g)}(t)\|
   \bh_v^{(g)}(t)]
\bigr).
\label{eq:triple_combined}
\end{equation}

\subsection{Elastic Weight Consolidation for Topology Evolution}

When the graph topology changes, for example through deployment events or infrastructure reconfiguration, a model retrained on the new topology may overwrite previously learned representations. Elastic weight consolidation~\cite{Kirkpatrick2017,Zenke2017} prevents this catastrophic forgetting by penalising changes to parameters that are important for earlier tasks. After learning on topology configuration~$\calT_i$ with optimal parameters~$\theta_i^*$, the loss for subsequent configuration~$\calT_{i+1}$ includes the regularisation term
\begin{equation}
\calL_{\mathrm{EWC}}(\theta)
= \calL_{\calT_{i+1}}(\theta)
  + \lambda_{\mathrm{ewc}}\sum_j F_j\,(\theta_j - \theta_{i,j}^*)^2,
\label{eq:ewc}
\end{equation}
where the diagonal Fisher information $F_j = \E[(\partial\log p(y|\bx,\theta_i^*)/\partial\theta_j)^2]$ quantifies the importance of parameter~$j$ for the previous configuration and $\lambda_{\mathrm{ewc}}$ controls the consolidation strength.

\subsection{Joint Multi-Level Classification Objective}

The model performs simultaneous classification at all three granularities through a multi-task loss,
\begin{equation}
\calL_{\mathrm{joint}}
= \lambda_s\calL_s + \lambda_r\calL_r + \lambda_n\calL_n
  + \lambda_{\mathrm{con}}\calL_{\mathrm{consist}},
\label{eq:triple_joint}
\end{equation}
where $\calL_g$ is the cross-entropy loss at granularity~$g$, $\lambda_g$ are balancing weights, and the consistency loss
\begin{equation}
\calL_{\mathrm{consist}}
= \sum_{(v_s,v_r)\in\calB_{sr}}
  \mathrm{KL}\!\bigl(p_s(y|v_s)\|p_r(y|v_r)\bigr)
+ \sum_{(v_s,v_n)\in\calB_{sn}}
  \mathrm{KL}\!\bigl(p_s(y|v_s)\|p_n(y|v_n)\bigr)
\label{eq:triple_consist}
\end{equation}
enforces agreement between the predictions of linked entities across abstraction levels.

\begin{algorithm}[htbp]
\caption{TripleE-TGNN Multi-Level Training}\label{alg:triple_e}
\KwInput{Hierarchical graphs $\{\calG_s,\calG_r,\calG_n\}$ with couplings $\calB_{sr},\calB_{sn},\calB_{rn}$, labels at each granularity, EWC strength $\lambda_{\mathrm{ewc}}$}
\KwOutput{Trained encoders $\mathrm{ENC}_s,\mathrm{ENC}_r,\mathrm{ENC}_n$ and cross-attention parameters}
Initialise all encoder and attention parameters\;
\For{each topology configuration $\calT_i$}{
  \For{epoch $= 1,\ldots,E$}{
    \For{each mini-batch}{
      \For{granularity $g \in \{s,r,n\}$}{
        Compute intra-granularity embeddings $\bh_v^{(g)}(t) \leftarrow \mathrm{ENC}_g(\bx_v^{(g)},\calN_v^{(g)},\Delta t)$\;
        Compute temporal attention $\tilde{\bh}_v^{(g)}$ and spatial attention $\hat{\bh}_v^{(g)}$\;
        Fuse: $\bh_v^{(g)} \leftarrow \mathrm{MLP}([\tilde{\bh}_v^{(g)}\|\hat{\bh}_v^{(g)}\|\bh_v^{(g)}])$\;
      }
      \For{each pair $(g,g')$ with $g\neq g'$}{
        Compute cross-granularity messages $\mathbf{m}_v^{(g\to g')}$ via bipartite attention\;
        Update: $\bh_v^{(g')} \leftarrow \mathrm{LayerNorm}(\bh_v^{(g')} + \mathbf{m}_v^{(g\to g')})$\;
      }
      Compute $\calL_{\mathrm{joint}} = \lambda_s\calL_s + \lambda_r\calL_r + \lambda_n\calL_n + \lambda_{\mathrm{con}}\calL_{\mathrm{consist}}$\;
      \If{$i > 1$}{
        Add EWC penalty: $\calL \leftarrow \calL_{\mathrm{joint}} + \lambda_{\mathrm{ewc}}\sum_j F_j(\theta_j-\theta_{i-1,j}^*)^2$\;
      }
      Backpropagate and update all parameters\;
    }
  }
  Store Fisher diagonal $\{F_j\}$ and optimal $\theta_i^*$ for EWC\;
}
\Return trained encoders and attention parameters
\end{algorithm}



%%=====================================================================
\section{Method M4: MambaShield --- State-Space Inference with Poisoning Resilience}
\label{sec:mamba}
%%=====================================================================

This section solves Sub-problem~4 by developing an efficient sequential model that satisfies the computational complexity constraint of Equation~\eqref{eq:ps_ssm} while maintaining robustness under training-time poisoning through progressive distillation and PAC-Bayesian certification. The method addresses Gap~4 (Adversarial $\times$ Efficiency; Non-stationary $\times$ Efficiency).

\subsection{Selective State-Space Formulation}

The selective state-space model~\cite{Gu2022,Gu2024} extends the classical linear system through input-dependent transitions. Given an input sequence $u(t)\in\R^{d_{\mathrm{in}}}$, the continuous-time dynamics are
\begin{equation}
\frac{d\bh(t)}{dt} = \bA(u(t))\,\bh(t) + \bB(u(t))\,u(t),
\qquad
y(t) = \bC\,\bh(t) + \bD\,u(t),
\label{eq:mamba_cts}
\end{equation}
where $\bA(u)=f_A(u;\theta_A)\in\R^{d_s\times d_s}$ and $\bB(u)=f_B(u;\theta_B)\in\R^{d_s\times d_{\mathrm{in}}}$ are computed by learned projections. Discretisation through zero-order hold at step size~$\Delta$ produces
\begin{equation}
\bh_k = \bar{\bA}_k\,\bh_{k-1} + \bar{\bB}_k\,u_k,
\qquad
\bar{\bA}_k = \exp(\bA(u_k)\Delta),
\quad
\bar{\bB}_k = \bA(u_k)^{-1}(\bar{\bA}_k - \mathbf{I})\bB(u_k),
\label{eq:mamba_disc}
\end{equation}
and the convolution kernel $\bar{\mathbf{K}}=(\bC\bar{\bA}^i\bar{\bB})_{i=0}^{L-1}$ is evaluated in $\calO(L\log L)$ time through the Fast Fourier Transform.

\subsection{Progressive Adversarial Robustness Distillation}

An ensemble of $M$ teacher models $\{f_{\theta_m}\}_{m=1}^M$, each trained on a disjoint data subset, provides robustness through knowledge distillation~\cite{Hinton2015,Xie2019}. The student model~$f_\phi$ learns from the teachers through the loss
\begin{equation}
\calL_{\mathrm{distill}}
= \sum_{m=1}^{M}w_m\;\mathrm{KL}\!\bigl(p_\phi(y|\bx)\|p_{\theta_m}(y|\bx)\bigr),
\label{eq:mamba_distill}
\end{equation}
where teacher weights~$w_m$ reflect validation performance. A curriculum schedule $\rho(t)=\min(\rho_0+\alpha t,\rho_{\max})$ progressively increases the fraction of corrupted samples in each mini-batch, and the total training loss interpolates between the task loss and the distillation loss,
\begin{equation}
\calL_{\mathrm{total}}
= (1-\lambda(t))\,\calL_{\mathrm{task}}
+ \lambda(t)\,\calL_{\mathrm{distill}},
\qquad
\lambda(t) = \min(\lambda_0+\beta t,\;1),
\label{eq:mamba_curriculum}
\end{equation}
so that teacher influence grows as the poisoning intensity increases.

\subsection{PAC-Bayesian Generalisation Bound}

\begin{theorem}[PAC-Bayesian Bound for Selective State-Space Models~\cite{McAllester1999,Catoni2007}]
\label{thm:pac_bayes}
For any $\delta\in(0,1)$, with probability at least $1-\delta$ over the draw of a training set~$S$ of size~$n$, for all posterior distributions~$Q$ over parameters,
\begin{equation}
\E_{\theta\sim Q}[R(\theta)]
\le \E_{\theta\sim Q}[\hat{R}_n(\theta)]
+ \sqrt{\frac{\mathrm{KL}(Q\|P)+\log(2\sqrt{n}/\delta)}{2n}},
\label{eq:pac_bayes}
\end{equation}
where $R(\theta)$ is the true risk, $\hat{R}_n(\theta)$ is the empirical risk, and $P$ is a data-independent prior.
\end{theorem}

\noindent\textit{Proof sketch.} The proof follows the standard PAC-Bayesian template of Catoni, applying a change of measure from $P$ to~$Q$ through the Donsker-Varadhan variational representation of the KL divergence, combined with a moment-generating-function argument for the empirical process. The bound is valid uniformly over all posteriors~$Q$, including those concentrated on a single parameter vector, and is therefore non-vacuous whenever the KL term is small relative to $n$. With Gaussian prior $P=\calN(\mathbf{0},\sigma_P^2\mathbf{I})$ and posterior $Q=\calN(\bm{\mu}_Q,\sigma_Q^2\mathbf{I})$, the KL term has the closed form
\begin{equation}
\mathrm{KL}(Q\|P)
= \frac{d}{2}\Bigl(
  \frac{\sigma_Q^2}{\sigma_P^2}
  + \frac{\|\bm{\mu}_Q\|^2}{\sigma_P^2}
  - 1
  - \log\frac{\sigma_Q^2}{\sigma_P^2}
\Bigr),
\label{eq:pac_kl_gauss}
\end{equation}
which is minimised by keeping the posterior close to the origin and its variance close to that of the prior. Full details appear in Appendix~A.\qed

\begin{algorithm}[htbp]
\caption{MambaShield: Training with Progressive Adversarial Robustness Distillation}\label{alg:mamba}
\KwInput{Training data $\calD$, teacher ensemble $\{f_{\theta_m}\}_{m=1}^M$, curriculum schedule $\rho_0,\alpha,\rho_{\max}$, PAC-Bayes prior $P$}
\KwOutput{Robust student model $f_\phi$ with PAC-Bayesian certificate}
Initialise student parameters $\phi_0$\;
Partition $\calD$ into $M$ disjoint subsets; train each teacher $f_{\theta_m}$ on its subset\;
Compute teacher weights $w_m \propto \mathrm{ValAcc}(f_{\theta_m})$\;
\For{iteration $t = 1,\ldots,T_{\mathrm{train}}$}{
  Set poisoning fraction $\rho(t) = \min(\rho_0 + \alpha t,\;\rho_{\max})$\;
  Set distillation weight $\lambda(t) = \min(\lambda_0 + \beta t,\;1)$\;
  Sample mini-batch; corrupt fraction $\rho(t)$ of labels\;
  \For{each sample $(\bx,y)$}{
    Discretise SSM: $\bar{\bA}_k = \exp(\bA(u_k)\Delta)$, $\bar{\bB}_k = \bA^{-1}(\bar{\bA}_k - \mathbf{I})\bB(u_k)$\;
    Run selective scan: $\bh_k = \bar{\bA}_k\bh_{k-1} + \bar{\bB}_k u_k$ for $k=1,\ldots,L$\;
    Compute student prediction $p_\phi(y|\bx)$\;
  }
  Compute $\calL_{\mathrm{task}} = \mathrm{CE}(p_\phi,y)$\;
  Compute $\calL_{\mathrm{distill}} = \sum_m w_m\,\mathrm{KL}(p_\phi\|p_{\theta_m})$\;
  Total: $\calL = (1-\lambda(t))\calL_{\mathrm{task}} + \lambda(t)\calL_{\mathrm{distill}} + \gamma\,\mathrm{KL}(Q_\phi\|P)$\;
  Update $\phi \leftarrow \phi - \eta\nabla_\phi\calL$\;
}
Evaluate PAC-Bayes bound: $R \le \hat{R}_n + \sqrt{(\mathrm{KL}(Q\|P)+\log(2\sqrt{n}/\delta))/(2n)}$\;
\Return $f_\phi$ and generalisation certificate
\end{algorithm}


%%=====================================================================
\section{Methods M5 and M6: Uncertainty-Calibrated Inference on Graphs}
\label{sec:stochastic}
%%=====================================================================

This section solves Sub-problem~6 by developing two complementary methods that achieve the sub-linear regret and uncertainty decomposition required by Equation~\eqref{eq:ps_regret}. Method M5 (Stochastic Transformer) provides variational Bayesian attention, and Method M6 (UC-HGP) provides hierarchical Gaussian process uncertainty at multiple temporal scales. Together they address Gap~6 (Adversarial $\times$ Accuracy; Non-stationary $\times$ Accuracy).

\subsection{Variational Bayesian Attention}

For an input sequence $\mathbf{X}=[\bx_1,\dots,\bx_n]\in\R^{n\times d}$, variational distributions~\cite{Blundell2015} are placed over the query, key, and value projection matrices,
\begin{equation}
q_{\phi_Q}(\bW_Q) = \calN(\bm{\mu}_Q,\mathrm{diag}(\bm{\sigma}_Q^2)),
\quad
q_{\phi_K}(\bW_K) = \calN(\bm{\mu}_K,\mathrm{diag}(\bm{\sigma}_K^2)),
\quad
q_{\phi_V}(\bW_V) = \calN(\bm{\mu}_V,\mathrm{diag}(\bm{\sigma}_V^2)).
\label{eq:stoch_var_attn}
\end{equation}
Samples are drawn through the reparameterisation trick, $\bW_Q=\bm{\mu}_Q+\bm{\sigma}_Q\odot\bm{\epsilon}_Q$ with $\bm{\epsilon}_Q\sim\calN(\mathbf{0},\mathbf{I})$, so that attention scores $A_{ij}=(\bx_i\bW_Q)^\top(\bx_j\bW_K)/\sqrt{d_k}$ become random variables whose uncertainty propagates to the output. Monte Carlo estimation with $M$ samples approximates the expected attention,
\begin{equation}
\E_{q_\phi}[\mathrm{Attention}(\mathbf{X})]
\approx \frac{1}{M}\sum_{m=1}^{M}
  \mathrm{Attention}\!\bigl(\mathbf{X};\bW_Q^{(m)},\bW_K^{(m)},\bW_V^{(m)}\bigr).
\label{eq:stoch_mc_attn}
\end{equation}

\subsection{Epistemic-Aleatoric Uncertainty Decomposition}

Total predictive uncertainty decomposes into epistemic uncertainty,
\begin{equation}
U_{\mathrm{epi}}
= \mathrm{Var}_{p(\theta|\calD)}\!\bigl[\E_{p(y|\bx,\theta)}[y]\bigr]
\approx \frac{1}{M}\sum_{m=1}^{M}(\hat{y}^{(m)}-\bar{y})^2,
\label{eq:stoch_epi}
\end{equation}
measuring the variance of the expected prediction across parameter samples, and aleatoric uncertainty,
\begin{equation}
U_{\mathrm{ale}}
= \E_{p(\theta|\calD)}\!\bigl[\mathrm{Var}_{p(y|\bx,\theta)}[y]\bigr]
\approx \frac{1}{M}\sum_{m=1}^{M}\hat{\sigma}^{(m)2},
\label{eq:stoch_ale}
\end{equation}
averaging the predicted variance over parameter samples, where $\hat{y}^{(m)}$ is the mean prediction and $\hat{\sigma}^{(m)2}$ the predicted variance from sample~$m$, and $\bar{y}=\frac{1}{M}\sum_m\hat{y}^{(m)}$.

\subsection{Stochastic Adversarial Training}

Rather than generating worst-case perturbations at fixed budgets, stochastic adversarial training draws perturbations from a learned distribution~\cite{Madry2018,Zhang2019trades,Shafahi2019},
\begin{equation}
\bdelta\sim q_\psi(\bdelta|\bx)
= \calN\!\bigl(\mu_\psi(\bx),\;\sigma_\psi(\bx)^2\mathbf{I}\bigr),
\label{eq:stoch_pert_dist}
\end{equation}
parameterised by neural networks $\mu_\psi,\sigma_\psi$ that learn spatially adaptive perturbation statistics. The training objective is the stochastic min-max problem
\begin{equation}
\min_\theta\;\max_\psi\;
\E_{(\bx,y)\sim\calD}\;
\E_{\bdelta\sim q_\psi(\bdelta|\bx)}\!
\bigl[\calL(f_\theta(\bx+\bdelta),y)\bigr],
\label{eq:stoch_adv_obj}
\end{equation}
solved by alternating gradient ascent on~$\psi$ and descent on~$\theta$.

\subsection{Multi-Objective Loss}

The full training loss combines five terms,
\begin{equation}
\calL_{\mathrm{total}}
= \lambda_1\calL_{\mathrm{task}}
+ \lambda_2\calL_{\mathrm{KL}}
+ \lambda_3\calL_{\mathrm{adv}}
+ \lambda_4\calL_{\mathrm{calib}}
+ \lambda_5\calL_{\mathrm{reg}},
\label{eq:stoch_total_loss}
\end{equation}
where $\calL_{\mathrm{task}}=\E_{(\bx,y)}[-\log p_\theta(y|\bx)]$ is the classification loss, $\calL_{\mathrm{KL}}~\cite{Blundell2015}=\mathrm{KL}(q_\phi(\theta)\|p(\theta))$ regularises the variational posterior, $\calL_{\mathrm{adv}}=\E_{(\bx,y)}\E_{\bdelta\sim q_\psi}[-\log p_\theta(y|\bx{+}\bdelta)]$ enforces adversarial robustness, $\calL_{\mathrm{calib}}=\mathrm{ECE}~\cite{Guo2017,Naeini2015}(\{(p_\theta(\hat{y}|\bx_i),\mathbf{1}[\hat{y}=y_i])\}_i)$ penalises miscalibration, and $\calL_{\mathrm{reg}}=\|\theta\|_2^2$ controls weight magnitude.

\begin{algorithm}[htbp]
\caption{Stochastic Transformer: Variational Bayesian Training with Adversarial Robustness}\label{alg:stoch_trans}
\KwInput{Training data $\calD$, prior $p(\theta)$, MC samples $M$, perturbation distribution parameters $\psi_0$, loss weights $\lambda_1,\ldots,\lambda_5$}
\KwOutput{Variational parameters $\phi^*$, perturbation parameters $\psi^*$}
Initialise variational parameters $\phi = \{\bm{\mu}_Q,\bm{\sigma}_Q,\bm{\mu}_K,\bm{\sigma}_K,\bm{\mu}_V,\bm{\sigma}_V\}$\;
Initialise perturbation network parameters $\psi$\;
\For{iteration $t = 1,\ldots,T$}{
  \For{each mini-batch $\{(\bx_i,y_i)\}$}{
    \textbf{Perturbation step:}
    Compute $\mu_\psi(\bx_i),\sigma_\psi(\bx_i)$ and sample $\bdelta_i\sim\calN(\mu_\psi(\bx_i),\sigma_\psi(\bx_i)^2\mathbf{I})$\;
    Update $\psi \leftarrow \psi + \alpha_\psi\nabla_\psi\frac{1}{M}\sum_{m=1}^M\calL(f_{\theta^{(m)}}(\bx_i+\bdelta_i),y_i)$\;
    \textbf{Model step:}
    \For{$m = 1,\ldots,M$}{
      Sample $\bW_Q^{(m)}=\bm{\mu}_Q+\bm{\sigma}_Q\odot\bm{\epsilon}^{(m)}$, similarly for $\bW_K^{(m)},\bW_V^{(m)}$\;
      Compute attention output and predictions $\hat{y}^{(m)},\hat{\sigma}^{(m)2}$\;
    }
    Compute $U_{\mathrm{epi}} = \frac{1}{M}\sum_m(\hat{y}^{(m)}-\bar{y})^2$, $U_{\mathrm{ale}} = \frac{1}{M}\sum_m\hat{\sigma}^{(m)2}$\;
    Compute $\calL_{\mathrm{total}} = \lambda_1\calL_{\mathrm{task}} + \lambda_2\calL_{\mathrm{KL}} + \lambda_3\calL_{\mathrm{adv}} + \lambda_4\calL_{\mathrm{calib}} + \lambda_5\calL_{\mathrm{reg}}$\;
    Update $\phi \leftarrow \phi - \alpha_\phi\nabla_\phi\calL_{\mathrm{total}}$\;
  }
}
\Return $\phi^*,\psi^*$ together with $U_{\mathrm{epi}},U_{\mathrm{ale}}$ for each test input
\end{algorithm}


%%=====================================================================
\section{Forward-Compatible Detection Under Protocol Regime Change}
\label{sec:pq_idps}
%%=====================================================================

This section addresses Sub-problem~5 (Gap~5: Non-stationary $\times$ Robustness) by constructing a detection architecture that maintains robustness when the data-generating protocols underlying the graph undergo structural change, providing the forward compatibility that the domain-specific foundation model (CyberSecLLM) requires.

\subsection{Protocol-Specific Feature Extraction}

When a new protocol regime is deployed, the observable features on graph edges, such as timing distributions, message sizes, and interaction counts, change in characteristic ways. For lattice-based protocols, for example, the timing of key-exchange operations exhibits variations that correlate with the computational structure of the underlying lattice operations. These variations are captured as a statistical feature vector,
\begin{equation}
\mathbf{f}_{\mathrm{lattice}}
= \bigl[\mu_t,\;\sigma_t,\;\mathrm{skew}(t),\;\mathrm{kurt}(t),\;
  \mathrm{pctl}(t),\;\mathrm{acf}(t)\bigr],
\label{eq:pq_feat_lattice}
\end{equation}
where $t=(t_1,\dots,t_n)$ denotes operational timing measurements, and the entries capture central tendency, dispersion, distribution shape, and temporal dependence. Analogous feature vectors $\mathbf{f}_{\mathrm{code}}$ and $\mathbf{f}_{\mathrm{hash}}$ are defined for code-based and hash-based protocol families, incorporating structural parameters such as key sizes, error weights, syndrome entropy, signature sizes, hash invocations, and Merkle tree depths.

\subsection{Unified Multi-Protocol Detection Architecture}

Protocol-specific encoders produce latent representations $\mathbf{z}_{\mathrm{lattice}}=\mathrm{ENC}_{\mathrm{lattice}}(\mathbf{f}_{\mathrm{lattice}})$, $\mathbf{z}_{\mathrm{code}}=\mathrm{ENC}_{\mathrm{code}}(\mathbf{f}_{\mathrm{code}})$, $\mathbf{z}_{\mathrm{hash}}=\mathrm{ENC}_{\mathrm{hash}}(\mathbf{f}_{\mathrm{hash}})$. Cross-attention fusion enables interaction between the three protocol representations,
\begin{equation}
\mathbf{z}_{\mathrm{fused}}
= \mathrm{MultiHeadAttention}\!\bigl(
  [\mathbf{z}_{\mathrm{lattice}},\;
   \mathbf{z}_{\mathrm{code}},\;
   \mathbf{z}_{\mathrm{hash}}]\bigr),
\label{eq:pq_fusion}
\end{equation}
and multi-task classification predicts both the category label and the target protocol,
\begin{equation}
p(\mathrm{class},\mathrm{protocol})
= \mathrm{Softmax}\!\bigl(\bW_{\mathrm{out}}\mathbf{z}_{\mathrm{fused}}
  + \mathbf{b}_{\mathrm{out}}\bigr).
\label{eq:pq_class}
\end{equation}

\subsection{Formal Security Reduction}

The robustness guarantee of the detection framework under the new protocol regime inherits from the hardness of the Learning With Errors problem~\cite{Regev2009,Peikert2016}. Given $\mathbf{A}\in\Z_q^{m\times n}$ and $\mathbf{b}=\mathbf{A}\mathbf{s}+\mathbf{e}\bmod q$ for secret~$\mathbf{s}$ and small error~$\mathbf{e}$, any polynomial-time adversary that distinguishes a legitimate from a manipulated protocol-negotiation sequence can be used to construct a solver for LWE, contradicting the assumed hardness. The formal reduction, presented in full in Appendix~A, establishes that the false-negative rate of the classifier is bounded by the advantage of the best LWE solver plus a statistical term that decreases with the number of training samples.

\begin{algorithm}[htbp]
\caption{Forward-Compatible Multi-Protocol Detection}\label{alg:pqidps}
\KwInput{Graph event stream with mixed protocol-regime edges, protocol-specific encoders, adversarial training budget $\epsilon$}
\KwOutput{Classification $(\mathrm{class},\mathrm{protocol})$ and LWE-based robustness certificate}
\For{each incoming graph event}{
  Identify protocol regime via protocol-negotiation sequence analysis\;
  \uIf{lattice-based (Kyber/Dilithium)}{
    Extract timing features $\mathbf{f}_{\mathrm{lattice}} = [\mu_t,\sigma_t,\mathrm{skew},\mathrm{kurt},\mathrm{pctl},\mathrm{acf}]$\;
    Encode: $\mathbf{z}_{\mathrm{lattice}} \leftarrow \mathrm{ENC}_{\mathrm{lattice}}(\mathbf{f}_{\mathrm{lattice}})$\;
  }\uElseIf{code-based (McEliece/Niederreiter)}{
    Extract structural features $\mathbf{f}_{\mathrm{code}}$; encode $\mathbf{z}_{\mathrm{code}} \leftarrow \mathrm{ENC}_{\mathrm{code}}(\mathbf{f}_{\mathrm{code}})$\;
  }\Else{
    Extract hash-based features $\mathbf{f}_{\mathrm{hash}}$; encode $\mathbf{z}_{\mathrm{hash}} \leftarrow \mathrm{ENC}_{\mathrm{hash}}(\mathbf{f}_{\mathrm{hash}})$\;
  }
  Fuse: $\mathbf{z}_{\mathrm{fused}} \leftarrow \mathrm{MultiHeadAttention}([\mathbf{z}_{\mathrm{lattice}},\mathbf{z}_{\mathrm{code}},\mathbf{z}_{\mathrm{hash}}])$\;
  Classify: $p(\mathrm{class},\mathrm{protocol}) \leftarrow \mathrm{Softmax}(\bW_{\mathrm{out}}\mathbf{z}_{\mathrm{fused}} + \mathbf{b}_{\mathrm{out}})$\;
}
\textbf{Training:} Alternate between standard forward passes and adversarial forward passes with PGD perturbations $\|\bdelta\|\le\epsilon$ applied to feature vectors\;
\textbf{Certification:} Verify LWE reduction (Appendix~A) and report robustness parameter $\lambda_{\mathrm{sec}}$\;
\Return classification and robustness certificate
\end{algorithm}


%%=====================================================================
\section{Method M7: Game-Theoretic Certification and Unified Framework}
\label{sec:game}
%%=====================================================================

This section solves Sub-problem~7 by developing the Stackelberg game formulation and the consistency coupling term $\sum_{i<j}\calL_{ij}(\theta_i,\theta_j)$ from the unified problem (Equation~\eqref{eq:ps_unified}). Together with the certification components, this constitutes Method M7 and addresses Gap~7 (all conditions $\times$ all requirements).

\subsection{Stackelberg Game Formulation}

The defender selects a model $f_\theta$ and a deployment configuration~$c$ comprising detection policy parameters. The adversary observes $(f_\theta,c)$ and chooses a perturbation strategy~$a$. Defender and adversary utilities are respectively
\begin{equation}
U_d(f_\theta,c,a)
= \E_{\bx\sim\calD(a)}\!\bigl[\mathbf{1}[f_\theta(\bx)=y]\bigr] - \mathrm{Cost}(c),
\label{eq:game_ud}
\end{equation}
\begin{equation}
U_a(f_\theta,c,a)
= \Prob[\text{misclassification}\mid f_\theta,c,a] - \mathrm{Cost}(a),
\label{eq:game_ua}
\end{equation}
where $\calD(a)$ is the data distribution induced by the adversary's strategy. The Stackelberg equilibrium is the solution to
\begin{equation}
(f_\theta^*,c^*)
= \arg\max_{f_\theta,c}\;U_d\!\bigl(f_\theta,c,\;a^*(f_\theta,c)\bigr),
\qquad
a^*(f_\theta,c)=\arg\max_a\;U_a(f_\theta,c,a).
\label{eq:game_stackel}
\end{equation}

\subsection{Nash Equilibrium Characterisation}

Under simultaneous play the mixed strategy profile $(\pi_d,\pi_a)$ is a Nash equilibrium when neither player improves utility by unilateral deviation,
\begin{equation}
\E_{\pi_a}\!\bigl[U_d(\pi_d,a)\bigr]
\ge \E_{\pi_a}\!\bigl[U_d(\pi_d',a)\bigr]\;\forall\,\pi_d',
\qquad
\E_{\pi_d}\!\bigl[U_a(f_\theta,\pi_a)\bigr]
\ge \E_{\pi_d}\!\bigl[U_a(f_\theta,\pi_a')\bigr]\;\forall\,\pi_a'.
\label{eq:game_nash}
\end{equation}
In the zero-sum special case $U_d+U_a=0$, the equilibrium coincides with the minimax value $\pi_d^*=\arg\max_{\pi_d}\min_a\E_{\pi_d}[U_d(f_\theta,a)]$.

\subsection{Online Learning with No-Regret Guarantees}

At round~$t$ the defender selects configuration $c_t\in\calC$ and observes reward $r_t=U_d(c_t,a_t)$. Cumulative regret after $T$ rounds is
\begin{equation}
R(T)
= \max_{c\in\calC}\sum_{t=1}^{T}U_d(c,a_t)
  - \sum_{t=1}^{T}U_d(c_t,a_t).
\label{eq:game_regret}
\end{equation}
The multiplicative weights update~\cite{AroraHK2012,CesaBianchi2006} maintains a distribution $p_t(c)\propto\exp\!\bigl(\eta\sum_{s=1}^{t-1}\mathbf{1}[c_s{=}c]\,r_s\bigr)$ with learning rate $\eta=\sqrt{\log|\calC|/T}$.

\begin{theorem}[Regret Bound]
\label{thm:regret}
The multiplicative weights algorithm with learning rate $\eta=\sqrt{\log|\calC|/T}$ achieves
\begin{equation}
R(T)\le 2\sqrt{T\log|\calC|},
\label{eq:game_regret_bound}
\end{equation}
so that the average per-round regret $R(T)/T=\calO(\sqrt{\log|\calC|/T})\to 0$ as $T\to\infty$.
\end{theorem}

\noindent\textit{Proof sketch.} The proof proceeds by bounding the potential function $\Phi_t=\log\sum_c\exp(\eta\sum_{s=1}^{t}r_s(c))$ from above and below. The upper bound uses the concavity of the logarithm and the definition of the algorithm; the lower bound isolates the best action in hindsight. Combining the two bounds and optimising over~$\eta$ yields the stated rate. The full argument appears in Appendix~A.\qed

\subsection{Uncertainty-Guided Strategy Selection}

Bayesian uncertainty from the stochastic transformer of Section~\ref{sec:stochastic} informs the defender's utility through an uncertainty-penalised payoff,
\begin{equation}
\tilde{U}_d(c,a)
= \begin{cases}
U_d(c,a) - \kappa\,U_{\mathrm{epi}} & \text{if}\;U_{\mathrm{epi}}<\tau,\\
-\infty & \text{otherwise},
\end{cases}
\label{eq:game_unc}
\end{equation}
where $\kappa$ discounts utility under uncertainty and threshold~$\tau$ triggers referral for expert review. This produces risk-sensitive policies that abstain from predictions when the model lacks sufficient confidence.

\begin{algorithm}[htbp]
\caption{Game-Theoretic Robustness with Uncertainty-Guided Online Learning}\label{alg:game}
\KwInput{Action space $\calC$, uncertainty threshold $\tau$, penalty $\kappa$, learning rate schedule}
\KwOutput{Defender policy $\pi_d^*$ with regret guarantee}
Initialise uniform distribution $p_1(c) = 1/|\calC|$ for all $c\in\calC$\;
Set $\eta = \sqrt{\log|\calC|/T}$\;
\For{round $t = 1,\ldots,T$}{
  Sample defender action $c_t \sim p_t$\;
  Observe adversary action $a_t$ and reward $r_t = U_d(c_t,a_t)$\;
  Query stochastic transformer for epistemic uncertainty $U_{\mathrm{epi}}$\;
  \uIf{$U_{\mathrm{epi}} \ge \tau$}{
    Set $\tilde{r}_t = -\infty$ (reject; route to expert review)\;
  }\Else{
    Set $\tilde{r}_t = r_t - \kappa\,U_{\mathrm{epi}}$\;
  }
  Update weights: $p_{t+1}(c) \propto p_t(c)\exp(\eta\,\tilde{r}_t\,\mathbf{1}[c_t=c])$ for all $c\in\calC$\;
  Normalise $p_{t+1}$\;
}
Compute empirical Nash equilibrium from play frequencies $\{(c_t,a_t)\}_{t=1}^T$\;
\Return $\pi_d^*$ and verified regret $R(T)\le 2\sqrt{T\log|\calC|}$
\end{algorithm}


%%=====================================================================
\subsection{Unified Certification Through Coupled Optimisation}
\label{sec:unified}
%%=====================================================================

The seven methods developed above (M1 through M6, plus the game-theoretic component of M7) share mathematical structures — adjoint sensitivity analysis, Wasserstein geometry, KL-divergence regularisation, and Lipschitz-bounded dynamics — that permit a unified treatment through coupled optimisation. This section completes Method M7 by developing the certification framework that integrates all component guarantees.

\subsection{Coupled Objective Function}

Let $\theta=\{\theta_{\mathrm{CT}},\theta_{\mathrm{TE}},\theta_{\mathrm{Fed}},\theta_{\mathrm{PQ}},\theta_{\mathrm{MS}},\theta_{\mathrm{ST}},\theta_{\mathrm{GT}}\}$ collect the parameters of all seven methods. The unified objective couples component losses with pairwise consistency terms,
\begin{equation}
\calL_{\mathrm{unified}}(\theta)
= \sum_{i=1}^{7}\lambda_i\,\calL_i(\theta_i)
+ \sum_{i<j}\lambda_{ij}\,\calL_{ij}(\theta_i,\theta_j),
\label{eq:unified_obj}
\end{equation}
where each $\calL_i$ is the method-specific loss from the corresponding section and the coupling loss
\begin{equation}
\calL_{ij}(\theta_i,\theta_j)
= \mathrm{KL}\!\bigl(p_{\theta_i}(y|\bx)\|p_{\theta_j}(y|\bx)\bigr)
+ \|\mathbf{z}_{\theta_i}(\bx)-\mathbf{z}_{\theta_j}(\bx)\|_2^2
\label{eq:unified_coupling}
\end{equation}
enforces consistency between the predictive distributions and the learned representations of methods~$i$ and~$j$.

\subsection{Convergence Analysis}

\begin{theorem}[Unified Framework Convergence]
\label{thm:unified}
Assume each component loss $\calL_i$ is $L_i$-smooth and each coupling loss $\calL_{ij}$ is $L_{ij}$-smooth. Under alternating gradient descent with learning rates $\eta_i\le 1/(L_i+\sum_j L_{ij})$, the unified optimisation converges at rate
\begin{equation}
\|\nabla\calL_{\mathrm{unified}}(\theta_T)\|^2
\le \frac{2\bigl(\calL_{\mathrm{unified}}(\theta_0)
  -\calL_{\mathrm{unified}}(\theta^*)\bigr)}{\eta_{\min}\,T},
\label{eq:unified_conv}
\end{equation}
where $\eta_{\min}=\min_i\eta_i$ and $\theta^*$ is a stationary point.
\end{theorem}

\noindent\textit{Proof sketch.} Each alternating gradient step on component~$i$ decreases the unified objective by at least $\eta_i\|\nabla_{\theta_i}\calL_{\mathrm{unified}}\|^2/2$ under the smoothness condition, because the coupling terms are smooth in each block of variables. Summing over all components and all $T$ rounds, dividing by $T$, and taking the minimum over learning rates yields the stated bound. The full derivation, including the treatment of cross-term gradients, appears in Appendix~A.\qed

\begin{algorithm}[htbp]
\caption{Unified Framework: Alternating Optimisation with Consistency Coupling}\label{alg:unified}
\KwInput{Component losses $\calL_1,\ldots,\calL_7$, coupling weights $\{\lambda_{ij}\}$, smoothness constants $\{L_i,L_{ij}\}$}
\KwOutput{Jointly optimised parameters $\theta^*=\{\theta_1^*,\ldots,\theta_7^*\}$}
Initialise all component parameters $\theta_1^{(0)},\ldots,\theta_7^{(0)}$\;
Set learning rates $\eta_i = 1/(L_i + \sum_j L_{ij})$ for each component\;
\For{round $t = 0,1,\ldots,T{-}1$}{
  \For{component $i = 1,\ldots,7$}{
    Compute component gradient $\mathbf{g}_i = \nabla_{\theta_i}\calL_i(\theta_i^{(t)})$\;
    Compute coupling gradients $\mathbf{g}_{ij} = \nabla_{\theta_i}\calL_{ij}(\theta_i^{(t)},\theta_j^{(t)})$ for all $j\neq i$\;
    Combined gradient: $\mathbf{g}_i^{\mathrm{total}} = \lambda_i\mathbf{g}_i + \sum_{j\neq i}\lambda_{ij}\mathbf{g}_{ij}$\;
    Update: $\theta_i^{(t+1)} \leftarrow \theta_i^{(t)} - \eta_i\,\mathbf{g}_i^{\mathrm{total}}$\;
  }
  Evaluate $\calL_{\mathrm{unified}}(\theta^{(t+1)})$ and check convergence criterion\;
}
\Return $\theta^* = \theta^{(T)}$
\end{algorithm}


%%=====================================================================
\section{Chapter Summary}
\label{sec:ch2_summary}
%%=====================================================================

This chapter has developed the complete mathematical foundations for seven novel methods, derived from a single unified problem statement (Equation~\eqref{eq:ps_unified}) through systematic decomposition into seven sub-problems. The unified problem was progressively constructed by incorporating adversarial robustness (Stage~2), continuous-time dynamics with Lipschitz certification (Stage~3), distributed learning with differential privacy (Stage~4), multi-level representation with anti-forgetting (Stage~5), computational efficiency through state-space models (Stage~6), and non-stationary adaptation with uncertainty quantification (Stage~7).

Method M1 (CT-TGNN) provides continuous-time graph dynamics with Lipschitz-certified robustness, solving Sub-problem~1 (Gap~1). Method M2 (FedLLM-API) achieves Byzantine-resilient federated optimisation with differential privacy and optimal-transport alignment, solving Sub-problem~2 (Gap~2). Method M3 (TripleE-TGNN) captures multi-level structure through triple embeddings with elastic weight consolidation, solving Sub-problem~3 (Gap~3). Method M4 (MambaShield) delivers efficient sequential inference with PAC-Bayesian generalisation certificates, solving Sub-problem~4 (Gap~4). The forward-compatible detection architecture addresses protocol regime changes, solving Sub-problem~5 (Gap~5). Methods M5/M6 (Stochastic Transformer and UC-HGP) provide calibrated Bayesian uncertainty under adversarial perturbation, solving Sub-problem~6 (Gap~6). Method M7 (Certification), comprising the game-theoretic framework and the unified optimisation coupling, ties all component guarantees together through consistency regularisation with a proven convergence rate (Theorem~\ref{thm:unified}), solving Sub-problem~7 (Gap~7).

Every method has been formulated at the level of mathematical operations on graph structures, without reference to any specific application domain. The theoretical guarantees --- convergence rates, robustness certificates, privacy budgets, regret bounds, and PAC-Bayesian generalisation bounds --- hold wherever these methods are deployed. The next chapter instantiates all methods in the domain of network-traffic graphs, describing the adaptation and alignment process through which the domain-independent mathematical framework is mapped onto the entities and constraints of cybersecurity and hybrid Network Intrusion Detection and Prevention Systems.

